\section*{Capítulo \ref{ch:b2:funcseq} --- Sequências e séries de
funções}

\begin{solution}{exo:b2:funcseq:1}
$f_n(x) = \frac{x}{1 + nx}$: limite \omterm{def:b2:funcseq:def}{pontual} $0$ em $\intcc{0}{1}$.
\omterm{def:b2:funcseq:def}{Uniformemente}: $f_n$ cresce em $\intcc{0}{1}$ (derivada
$\frac{1}{(1+nx)^2} > 0$), logo $\norm{f_n}_\infty = f_n(1) =
\frac{1}{1+n} \to 0$: \emph{uniforme} em $\intcc{0}{1}$.

$g_n(x) = nx\,\eu^{-nx^2}$: limite \omterm{def:b2:funcseq:def}{pontual} $0$ (a exponencial
vence). Sup: $g_n' = n\eu^{-nx^2}(1 - 2nx^2)$ anula-se em $x_n =
\frac{1}{\sqrt{2n}}$, onde $g_n(x_n) = \sqrt{\frac n2}\,\eu^{-1/2}
\to \infty$: não uniforme em $\intcc{0}{1}$ --- mas uniforme em
$\intco{\delta}{1}$, pois aí $g_n(x) \leq n\,\eu^{-n\delta^2}
\to 0$.

$h_n(x) = x^n(1 - x^n)$: limite \omterm{def:b2:funcseq:def}{pontual} $0$ em $\intcc{0}{1}$
(os dois fatores; em $x = 1$, $h_n = 0$). Sup: com $u = x^n \in
\intcc{0}{1}$, $u(1-u) \leq \frac14$ atingido em $u = \frac12$,
isto é, $x = 2^{-1/n} \in \intoo{0}{1}$: $\norm{h_n}_\infty =
\frac14 \not\to 0$: não uniforme em $\intcc{0}{1}$; uniforme em
$\intcc{0}{a}$ ($\sup \leq a^n \to 0$).
\end{solution}

\begin{solution}{exo:b2:funcseq:2}
$\int_0^1 nx\,\eu^{-nx^2}\dd x = \bigl[-\tfrac12
\eu^{-nx^2}\bigr]_0^1 = \frac{1 - \eu^{-n}}{2} \to \frac12$, enquanto
$\int_0^1 \lim g_n = 0$. Não há contradição:
o~\cref{thm:b2:funcseq:integration} exige convergência \emph{uniforme}
no segmento, o que aqui falha (a corcova de altura
$\sim\sqrt n$ desliza rumo a $0$).
\end{solution}

\begin{solution}{exo:b2:funcseq:3}
\omterm{def:b2:funcseq:series}{Convergência normal} em $\intcc{-1}{1}$: $\norm{x^n/n^2}_\infty =
\frac{1}{n^2}$, \omterm{def:b2:series:summable}{somável}: $S$ é \omterm{def:b2:metric:continuity}{contínua} aí
(\cref{thm:b2:funcseq:seriestransfer}).

Derivada: a série derivada $\sum \frac{x^{n-1}}{n}$ converge
\omterm{def:b2:funcseq:series}{normalmente} em todo $\intcc{-a}{a}$, $a < 1$ ($\sup =
\frac{a^{n-1}}{n}$): $S$ é $C^1$ em $\intoo{-1}{1}$ com
\[
S'(x) = \sum_{n\geq1} \frac{x^{n-1}}{n} = \frac1x \sum_{n\geq1}
\frac{x^n}{n} = -\frac{\ln(1 - x)}{x}
\qquad (0 < \abs x < 1),
\]
sendo a última identidade a série logarítmica do primeiro ano (redemonstrada
honestamente no~\cref{ch:b2:powerseries}).
\end{solution}

\begin{solution}{exo:b2:funcseq:4}
Para $x > 0$ fixo a série é alternada com $\frac{1}{n + x}
\downarrow 0$: \omterm{def:b2:funcseq:def}{convergência pontual}, e a cota do resto
$\abs{R_N(x)} \leq \frac{1}{N + 1 + x} \leq \frac{1}{N+1}$ é
\emph{uniforme} em $\intco{\delta}{\infty}$ (de fato, em
$\intoo{0}{\infty}$): \omterm{def:b2:funcseq:def}{convergência uniforme}. (Não normal:
$\norm{\frac{(-1)^n}{n+x}}_\infty = \frac{1}{n + \delta}$,
divergente.) A \omterm{def:b2:metric:continuity}{continuidade} decorre do
\cref{thm:b2:funcseq:seriestransfer}.

Equação funcional: reindexe $F(x + 1)$ com $m = n + 1$:
\[
F(x+1) = \sum_{n\geq0} \frac{(-1)^n}{n + 1 + x}
= \sum_{m\geq1}\frac{(-1)^{m-1}}{m+x} ,
\]
de modo que, isolando o termo $m = 0$ de $F(x)$,
\[
F(x) + F(x+1)
= \frac{1}{x} + \sum_{m\geq1}
\frac{(-1)^m + (-1)^{m-1}}{m+x} = \frac1x .
\]
\end{solution}

\begin{solution}{exo:b2:funcseq:5}
Ponha $g_n = f_n - f \geq 0$ (decrescente em $n$, por hipótese; o
limite é $0$ \omterm{def:b2:funcseq:def}{pontualmente}); cada $g_n$ é \omterm{def:b2:metric:continuity}{contínua}. Fixe $\varepsilon
> 0$ e ponha $U_n = \{x : g_n(x) < \varepsilon\}$: \omterm{def:b2:metric:topology}{aberto} (imagem inversa
de um \omterm{def:b2:metric:topology}{aberto}), crescente ($g_{n+1} \leq g_n$) e cobrindo $K$
(\omterm{def:b2:funcseq:def}{convergência pontual}). Por Borel--Lebesgue
(\cref{thm:b2:metric:borellebesgue}), um número finito de $U_{n_1}
\subseteq \dots \subseteq U_{n_k}$ cobre $K$: logo $K = U_{n_k}$,
isto é, $\norm{g_{n_k}}_\infty \leq \varepsilon$ e, pela monotonicidade,
$\norm{g_n}_\infty \leq \varepsilon$ para todo $n \geq n_k$: \omterm{def:b2:funcseq:def}{convergência
uniforme}. (A monotonicidade é essencial: as corcovas deslizantes do
\cref{exo:b2:funcseq:2} convergem \omterm{def:b2:funcseq:def}{pontualmente} num \omterm{def:b2:metric:compact}{compacto} sem
uniformidade.)
\end{solution}

\begin{solution}{exo:b2:funcseq:6}
Substitua $u = nx$:
\[
\int_0^1 \frac{n f(x)}{1 + n^2x^2}\dd x
= \int_0^n \frac{f(u/n)}{1 + u^2}\,\dd u .
\]
Os integrandos $h_n(u) = \frac{f(u/n)}{1+u^2}\mathbf{1}_{u \leq n}$
convergem \omterm{def:b2:funcseq:def}{pontualmente} para $\frac{f(0)}{1+u^2}$ (\omterm{def:b2:metric:continuity}{continuidade} de $f$ em
$0$) e são dominados por $\frac{\norm f_\infty}{1 + u^2}$,
integrável em $\intco{0}{\infty}$: a convergência dominada
(\cref{thm:b2:integration:dominated}) dá o limite
\[
\int_0^\infty \frac{f(0)}{1 + u^2}\dd u = \frac{\pi}{2} f(0) .
\]
(Os núcleos se concentram em $0$: uma identidade aproximada.)
\end{solution}

\begin{solution}{exo:b2:funcseq:7}
Para $f(x) = x^2$, use a segunda família de identidades binomiais da
demonstração: $\sum_k k^2 p_k = n(n-1)x^2 + nx$. Logo
\[
B_n(f)(x) = \sum_k \frac{k^2}{n^2}\,p_k
= \frac{n(n-1)x^2 + nx}{n^2}
= x^2 + \frac{x(1 - x)}{n} :
\]
$B_n(f) - f = \frac{x(1-x)}{n}$, de \omterm{def:b2:nvs:norm}{norma} do sup $\frac{1}{4n} =
O\bigl(\frac1n\bigr)$, como previsto.
\end{solution}

\begin{solution}{exo:b2:funcseq:8}
Para $\varepsilon = 1$ existe $N$ com $\norm{P_n - P_m}_{\infty,
\R} \leq 1$ para $m, n \geq N$. Um polinômio limitado em $\R$ é
constante (um não constante tende a $\pm\infty$): $P_n - P_m =
c_{n,m}$, constantes. Logo, para $n \geq N$: $P_n = P_N + c_n$ com
$c_n = P_n(0) - P_N(0)$ convergente (\omterm{def:b2:funcseq:def}{convergência pontual} em $0$).
Portanto $f = \lim P_n = P_N + \lim c_n$: um polinômio.
\end{solution}

\begin{solution}{exo:b2:funcseq:9}
(a) $\norm{(3/4)^n\varphi(4^n\cdot)}_\infty = \frac12
(3/4)^n$: \omterm{def:b2:funcseq:series}{convergência normal}, logo $W$ é \omterm{def:b2:metric:continuity}{contínua}
(\cref{thm:b2:funcseq:seriestransfer}).

(b) Fixe $x$, $m$; escolha o sinal de $h_m = \pm\frac12 4^{-m}$ de modo
que o segmento $\intcc{4^mx}{4^m(x + h_m)}$ (de comprimento $\frac12$)
não contenha semi-inteiro algum, tornando $\varphi$ afim de inclinação $\pm1$
nele (possível: um intervalo de comprimento $\frac12$ encontra no máximo um
ponto semi-inteiro; escolha o lado que o evita).

Para $n > m$: $4^n h_m = \pm\frac12 4^{n-m}$ é inteiro, e
$\varphi$ é $1$-periódica: o $n$-ésimo termo da diferença
se anula.

Para $n = m$: $\abs{\varphi(4^m x + 4^m h_m) - \varphi(4^m x)} =
\abs{4^m h_m} = \frac12$ ($\varphi$ afim de inclinação $\pm 1$ no
segmento), de modo que o termo contribui exatamente com $(3/4)^m \cdot
\frac{1/2}{\abs{h_m}} = (3/4)^m\,4^m = 3^m$ no quociente.

Para $n < m$: a $1$-lipschitziana $\varphi$ dá
$\bigl|(3/4)^n\bigl(\varphi(4^nx + 4^nh_m) -
\varphi(4^nx)\bigr)\bigr| \leq (3/4)^n 4^n\abs{h_m} =
3^n\abs{h_m}$: cada um contribui no máximo com $3^n$ para o quociente.

Logo
\[
\Bigl|\frac{W(x + h_m) - W(x)}{h_m}\Bigr|
\geq 3^m - \sum_{n=0}^{m-1} 3^n
= 3^m - \frac{3^m - 1}{2} = \frac{3^m + 1}{2}
\longrightarrow \infty .
\]
Se $W$ fosse derivável em $x$, todo quociente de diferenças ao longo de
$h_m \to 0$ convergiria para $W'(x)$: contradição. $W$ é
\omterm{def:b2:metric:continuity}{contínua} em toda parte e derivável em nenhuma.
\end{solution}

\begin{solution}{exo:b2:funcseq:10}
\omterm{def:b2:funcseq:def}{Pontualmente}: para $x \in \intco{0}{1}$ a série é geométrica de
razão $-x$,
\[
\sum_{n\geq0}(-1)^n x^n(1-x) = \frac{1-x}{1+x},
\]
e em $x = 1$ todo termo se anula: soma $0 = \frac{1-1}{2}$,
coerente --- a soma é \omterm{def:b2:metric:continuity}{contínua} em $\intcc{0}{1}$.
Uniformidade: o resto é uma cauda geométrica,
\[
\abs{R_N(x)} = \frac{x^{N+1}(1-x)}{1+x} \leq x^{N+1}(1 - x)
\leq \max_{\intcc01} t^{N+1}(1-t)
= \frac{1}{N+2}\Bigl(\frac{N+1}{N+2}\Bigr)^{\!N+1}
\leq \frac{1}{N+2} \to 0 ,
\]
uniformemente em $x$. Não normal: $\norm{u_n}_\infty =
\max x^n(1-x) = \frac{1}{n+1}\bigl(\frac{n}{n+1}\bigr)^n \sim
\frac{1}{\eu\,n}$, e $\sum \frac1{\eu n}$ diverge. Contraste:
$\sum x^n(1-x)$ tem somas parciais $1 - x^{N+1}$, convergindo
\omterm{def:b2:funcseq:def}{pontualmente} para a \emph{descontínua} $\mathbf 1_{\intco01}$: pelo
\cref{thm:b2:funcseq:continuity}, essa convergência não pode ser
uniforme em $\intcc{0}{1}$.
\end{solution}

\begin{solution}{exo:b2:funcseq:11}
O limite $f$ é \omterm{def:b2:metric:continuity}{contínuo} (\cref{thm:b2:funcseq:continuity}).
Então
\[
\abs{f_n(x_n) - f(x)}
\leq \abs{f_n(x_n) - f(x_n)} + \abs{f(x_n) - f(x)}
\leq \norm{f_n - f}_\infty + \abs{f(x_n) - f(x)} ,
\]
e os dois termos tendem a $0$ (\omterm{def:b2:funcseq:def}{convergência uniforme}; \omterm{def:b2:metric:continuity}{continuidade} de
$f$ em $x$). Contraexemplo sob mera \omterm{def:b2:funcseq:def}{convergência pontual}:
$g_n(x) = nx\,\eu^{-nx^2} \to 0$ \omterm{def:b2:funcseq:def}{pontualmente} em $\intcc{0}{1}$
com $g_n$ e o limite \omterm{def:b2:metric:continuity}{contínuo}, e no entanto, em $x_n =
\frac{1}{\sqrt{2n}} \to 0$:
\[
g_n(x_n) = \sqrt{\frac n2}\,\eu^{-1/2} \longrightarrow +\infty
\neq 0 = f(0) .
\]
\end{solution}

\begin{solution}{exo:b2:funcseq:12}
\begin{enumerate}
  \item Indução. $n = 1$ é a definição. Suponha a fórmula
        para $n$ e ponha $g(x) = \int_0^x \frac{(x-t)^n}{n!}
        f(t)\dd t$. Para um integrando \omterm{def:b2:metric:continuity}{contínuo} em $(x,t)$ e
        $C^1$ em $x$, a \omterm{thm:b2:integration:continuity}{integral com parâmetro} de limite variável
        deriva-se como
        \[
        g'(x) = \frac{(x-x)^n}{n!}f(x)
        + \int_0^x \frac{(x-t)^{n-1}}{(n-1)!}f(t)\dd t
        = T^nf(x)
        \]
        (separe $g(x+h) - g(x)$ na faixa
        $\int_x^{x+h}$, que é $O(h\cdot\sup)$ com o
        integrando anulando-se em $t = x$ como $h^n$, e na
        integral fixa do incremento em $x$, tratada pela desigualdade do valor
        médio e pela \omterm{def:b2:metric:continuity}{continuidade}). Além disso, $(T^{n+1}f)' = T^nf$
        (teorema fundamental do cálculo) e $g(0) =
        T^{n+1}f(0) = 0$: duas primitivas de $T^nf$ que se anulam em
        $0$ coincidem, logo $T^{n+1}f = g$. A cota:
        \[
        \abs{T^nf(x)} \leq \norm f_\infty
        \int_0^x \frac{(x-t)^{n-1}}{(n-1)!}\dd t
        = \norm f_\infty\,\frac{x^n}{n!}
        \leq \frac{\norm f_\infty}{n!} .
        \]
  \item $\sum_n \norm{T^nf}_\infty \leq \eu\,\norm f_\infty$:
        \omterm{def:b2:funcseq:series}{convergência normal}, logo uniforme; $S$ é \omterm{def:b2:metric:continuity}{contínua}.
        As somas parciais satisfazem $S_N = f + T S_{N-1}$, e $T$
        é $1$-lipschitziana para $\norm\cdot_\infty$
        ($\abs{Tg(x)} \leq x\norm g_\infty$): fazendo $N \to
        \infty$ nos dois lados obtém-se $S = f + TS$.
  \item Ponha $V(x) = f(x) + \eu^x\int_0^x \eu^{-t}f(t)\dd t$.
        Então $V - f$ é $C^1$ com $(V-f)'(x) =
        \eu^x\int_0^x\eu^{-t}f + f(x) = V(x)$, e $(TV)' = V$
        com $(V - f)(0) = TV(0) = 0$: logo $V - f = TV$, isto é,
        $V$ resolve a equação. Unicidade: se $S_1, S_2$ são
        soluções \omterm{def:b2:metric:continuity}{contínuas}, $D = S_1 - S_2$ satisfaz $D =
        TD$, logo $D = T^nD$ para todo $n$ e $\norm D_\infty
        \leq \frac{\norm D_\infty}{n!} \to 0$: $D = 0$.
        Portanto
        \[
        \sum_{n\geq0} T^nf(x) = f(x) +
        \int_0^x \eu^{x-t}f(t)\,\dd t .
        \]
        (A série $\sum T^n$ é uma série geométrica de
        operadores: um primeiro gostinho do resolvente
        $(\mathrm{Id} - T)^{-1}$, desenvolvido no volume do terceiro ano
        de graduação.)
\end{enumerate}
\end{solution}

\begin{solution}{pb:b2:funcseq:1}
\textbf{1.} A linearidade é clara pela fórmula. Positividade:
os pesos $p_k(x) \geq 0$, de modo que $f \geq 0$ força $B_nf \geq
0$; a monotonicidade segue aplicando isso a $g - f$. Cota: $\pm f \leq
\norm f_\infty$ dá $\pm B_nf \leq \norm f_\infty B_ne_0 =
\norm f_\infty$. Extremidades: $p_k(0) = \mathbf 1_{k=0}$ e
$p_k(1) = \mathbf 1_{k=n}$, logo $B_nf(0) = f(0)$, $B_nf(1) =
f(1)$.

\textbf{2.} Derive $(x+y)^n = \sum_k\binom nk x^ky^{n-k}$
em $x$, multiplique por $x$ e ponha $y = 1 - x$:
\[
nx = \sum_k k\,p_k(x) ;
\]
duas vezes, multiplicando por $x^2$: $n(n-1)x^2 = \sum_k k(k-1)p_k(x)$.
Logo $B_ne_0 = 1$ (binômio de Newton), $B_ne_1(x) =
\frac{nx}{n} = x$ e
\[
B_ne_2(x) = \frac{\sum_k k^2p_k}{n^2}
= \frac{n(n-1)x^2 + nx}{n^2}
= x^2 + \frac{x(1-x)}{n} .
\]

\textbf{3.} Desenvolva:
\[
\sum_k\Bigl(\frac kn - x\Bigr)^{\!2} p_k
= B_ne_2(x) - 2x\,B_ne_1(x) + x^2
= \frac{x(1-x)}{n} .
\]
Cauchy--Schwarz com a decomposição $\abs{k/n - x}\sqrt{p_k}
\cdot \sqrt{p_k}$:
\[
\sum_k\Bigl|\frac kn - x\Bigr| p_k
\leq \Bigl(\sum_k\Bigl(\frac kn - x\Bigr)^2
p_k\Bigr)^{\!1/2}
= \sqrt{\frac{x(1-x)}{n}} \leq \frac{1}{2\sqrt n},
\]
usando $x(1-x) \leq \frac14$.

\textbf{4.} Os pesos $p_k(x)$ são não negativos, somam $1$
e têm baricentro $\sum_k \frac kn p_k(x) = x$ (questão 2). A
desigualdade de Jensen finita para a convexa $f$ (indução a partir da
definição com dois pontos, volume do primeiro ano de graduação) dá
\[
f(x) = f\Bigl(\sum_k \frac kn\,p_k\Bigr)
\leq \sum_k f\Bigl(\frac kn\Bigr)p_k = B_nf(x) .
\]

\textbf{5.} Em $\{k : \abs{k/n - x} > \delta\}$ tem-se
$\bigl(\frac{k/n - x}{\delta}\bigr)^2 > 1$, logo
\[
\sum_{\abs{k/n-x}>\delta} p_k
\leq \frac{1}{\delta^2}\sum_k\Bigl(\frac kn -
x\Bigr)^2p_k
= \frac{x(1-x)}{n\delta^2} \leq \frac{1}{4n\delta^2} .
\]
Leitura: $p_k(x)$ é a lei de uma frequência amostral $S_n/n$ de
$n$ lançamentos de moeda de viés $x$; sua média é $x$, sua variância
$\frac{x(1-x)}n \to 0$, e a fórmula exibida é a desigualdade de Chebyshev:
a massa se concentra em $x$, de modo que promediar $f$
contra ela reproduz $f(x)$ no limite
(o~\cref{ch:b2:genfun} torna o vocabulário oficial).

\textbf{6.} $\omega \leq 2M < \infty$; a monotonicidade é clara
(sup sobre um conjunto maior). Heine: $f$ \omterm{def:b2:metric:continuity}{contínua} num \omterm{def:b2:metric:compact}{compacto} é
uniformemente contínua, o que diz exatamente que $\omega(\delta) \to 0$
quando $\delta \to 0^+$. Subaditividade: se $\abs{s - t} \leq
\delta_1 + \delta_2$, o ponto $u$ do segmento $\intcc st$
a distância $\min(\delta_1, \abs{s-t})$ de $s$ satisfaz
$\abs{s-u} \leq \delta_1$, $\abs{u-t} \leq \delta_2$ e
$\abs{f(s)-f(t)} \leq \abs{f(s)-f(u)} + \abs{f(u)-f(t)}$.
Iterando, $\omega(p\delta) \leq p\,\omega(\delta)$ para $p \in
\N^*$; para $\lambda > 0$, com $p = \lceil\lambda\rceil \leq 1
+ \lambda$: $\omega(\lambda\delta) \leq \omega(p\delta) \leq
p\,\omega(\delta) \leq (1+\lambda)\omega(\delta)$.

\textbf{7.} Como $\sum p_k = 1$:
\[
\abs{B_nf(x) - f(x)}
= \Bigl|\sum_k\bigl(f(k/n) - f(x)\bigr)p_k\Bigr|
\leq \sum_k\omega\bigl(\abs{k/n - x}\bigr)p_k .
\]
Para cada $k$, a questão 6 com $\lambda = \abs{k/n - x}/\delta$
dá $\omega(\abs{k/n-x}) \leq \bigl(1 +
\frac{\abs{k/n-x}}\delta\bigr)\omega(\delta)$; somando contra
os $p_k$ obtém-se a estimativa mestra.

\textbf{8.} Insira a cota da questão 3:
\[
\abs{B_nf(x) - f(x)} \leq \Bigl(1 +
\frac{1}{2\delta\sqrt n}\Bigr)\omega(\delta),
\]
uniformemente em $x$; com $\delta = n^{-1/2}$ o parêntese vale
$\frac32$: $\norm{B_nf - f}_\infty \leq
\frac32\omega(n^{-1/2}) \to 0$ pela questão 6 (Heine). Isso
redemonstra o~\cref{thm:b2:funcseq:weierstrass} com uma taxa.

\textbf{9.} $L$-lipschitziana significa $\omega(\delta) \leq L\delta$:
taxa $\frac{3L}{2\sqrt n}$. $\alpha$-hölderiana significa
$\omega(\delta) \leq C\delta^\alpha$: taxa $\frac{3C}{2}
n^{-\alpha/2}$.

\textbf{10.} Usando $k\binom{2m}k = 2m\binom{2m-1}{k-1}$:
\[
\sum_{k=m+1}^{2m}k\binom{2m}k
= 2m\sum_{j=m}^{2m-1}\binom{2m-1}{j}
= 2m\cdot 2^{2m-2},
\]
porque $j \mapsto 2m-1-j$ leva bijetivamente $\{m,\dots,2m-1\}$ em
$\{0,\dots,m-1\}$, de modo que a soma é metade de $2^{2m-1}$. Além disso,
$\sum_{k=m+1}^{2m}\binom{2m}k = \frac{2^{2m} -
\binom{2m}m}{2}$ (mesma simetria). Logo
\[
\sum_{k=m+1}^{2m}(k-m)\binom{2m}k
= m\,2^{2m-1} - m\,\frac{2^{2m} - \binom{2m}m}{2}
= \frac m2\binom{2m}m .
\]
A substituição $k \mapsto 2m-k$ leva os termos com $k < m$
nos com $k > m$ (binomiais iguais, $\abs{k-m}$ iguais):
a soma absoluta é o dobro da soma de um lado, $m\binom{2m}m$.

\textbf{11.} Em $x = \frac12$, $p_k(\tfrac12) =
\binom{2m}k2^{-2m}$ e $f(\tfrac12) = 0$:
\[
B_{2m}f\Bigl(\frac12\Bigr)
= \sum_k\Bigl|\frac{k}{2m} - \frac12\Bigr|
\binom{2m}k 2^{-2m}
= \frac{2^{-2m}}{2m}\,m\binom{2m}m
= \frac{\binom{2m}m}{2\cdot4^m}
\sim \frac{1}{2\sqrt{\pi m}}
\]
pelo~\cref{ex:b2:comparison:centralbinomial}. Como aqui $\omega_f
(\delta) = \delta$ (a função é $1$-lipschitziana e a
cota é atingida), a questão 8 prevê no máximo
$\frac32(2m)^{-1/2}$: o erro verdadeiro $\frac{1}{2\sqrt{\pi m}}$
tem exatamente a ordem $n^{-1/2}$ --- a taxa é ótima a menos da
constante.

\textbf{12.} $f \leq g$ dá $g - f \geq 0$, logo $L(g-f) \geq
0$, isto é, $Lf \leq Lg$. De $-\abs f \leq f \leq \abs f$:
$-L\abs f \leq Lf \leq L\abs f$, isto é, $\abs{Lf} \leq
L\abs f$.

\textbf{13.} Por Heine escolha $\delta$ com $\abs{f(s)-f(x)}
\leq \varepsilon$ sempre que $\abs{s-x} \leq \delta$. Se
$\abs{s - x} > \delta$, então $\frac{(s-x)^2}{\delta^2} > 1$ e
$\abs{f(s)-f(x)} \leq 2M \leq \frac{2M}{\delta^2}(s-x)^2$. Nos
dois casos a cota afirmada vale.

\textbf{14.} Fixe $x$; a questão 13 diz, como funções de $s$:
\[
-\varepsilon e_0 - \frac{2M}{\delta^2}q_x
\;\leq\; f - f(x)e_0
\;\leq\; \varepsilon e_0 + \frac{2M}{\delta^2}q_x,
\qquad q_x = e_2 - 2x\,e_1 + x^2e_0 .
\]
Aplique o $L_n$ monótono linear (questão 12) e avalie em
$x$:
\[
\abs{L_nf(x) - f(x)L_ne_0(x)}
\leq \varepsilon L_ne_0(x)
+ \frac{2M}{\delta^2}\bigl(L_ne_2(x) - 2xL_ne_1(x)
+ x^2L_ne_0(x)\bigr) .
\]

\textbf{15.} Escreva $\alpha_j = L_ne_j - e_j$, de modo que
$\norm{\alpha_j}_\infty \to 0$. Como $e_2(x) - 2xe_1(x) +
x^2e_0(x) = 0$:
\[
L_ne_2(x) - 2xL_ne_1(x) + x^2L_ne_0(x)
= \alpha_2(x) - 2x\,\alpha_1(x) + x^2\alpha_0(x),
\]
de \omterm{def:b2:nvs:norm}{norma} do sup no máximo $\norm{\alpha_2} + 2\norm{\alpha_1} +
\norm{\alpha_0} \to 0$. Além disso, $L_ne_0 \to e_0$ \omterm{def:b2:funcseq:def}{uniformemente}, logo
$L_ne_0 \leq 2$ para $n$ grande, e $\abs{f(x)}\abs{L_ne_0(x) -
1} \leq M\norm{\alpha_0} \to 0$. Juntando com a questão 14:
para $n$ grande, uniformemente em $x$,
\[
\abs{L_nf(x) - f(x)} \leq 2\varepsilon +
\frac{2M}{\delta^2}\,o(1) + M\,o(1) \leq 3\varepsilon :
\]
$L_nf \to f$ \omterm{def:b2:funcseq:def}{uniformemente} --- o teorema de Korovkin.

\textbf{16.} $B_ne_0 = e_0$ e $B_ne_1 = e_1$ exatamente, e
$\norm{B_ne_2 - e_2}_\infty = \max_x\frac{x(1-x)}{n} =
\frac{1}{4n} \to 0$ (questão 2): Korovkin se aplica, e
Weierstrass segue pela terceira vez.

\textbf{17.} $I_nf$ é linear em $f$ (os valores nodais o são) e, em
cada célula, a interpolante afim de valores nodais não negativos é
não negativa: positivo. $I_ne_0 = e_0$ e $I_ne_1 = e_1$
porque uma função afim é igual à própria interpolante. Numa
célula $\intcc ab$ ($b - a = \frac1n$), a interpolante afim de
$e_2$ é $L(t) = (a+b)t - ab$, e
\[
L(t) - t^2 = (t-a)(b-t) \in
\intcc{0}{\tfrac{(b-a)^2}{4}} ,
\]
com o máximo no ponto médio: $\norm{I_ne_2 - e_2}_\infty =
\frac{1}{4n^2} \to 0$. Korovkin: $I_nf \to f$ \omterm{def:b2:funcseq:def}{uniformemente} para
toda $f$ \omterm{def:b2:metric:continuity}{contínua} --- aproximação poligonal, sem
necessidade de mais nenhuma estimativa.

\textbf{18.} Dados $f$ e $\varepsilon$: Weierstrass fornece
um polinômio $P = \sum_{j=0}^d a_jx^j$ com $\norm{f -
P}_\infty \leq \frac\varepsilon2$; substituir cada $a_j$ por um
racional $b_j$ com $\abs{a_j - b_j} \leq
\frac{\varepsilon}{2(d+1)}$ move a \omterm{def:b2:nvs:norm}{norma} do sup em
$\intcc{0}{1}$ em no máximo $\frac\varepsilon2$. O conjunto dos
polinômios de coeficientes racionais é uma união \omterm{def:b2:structures:countable}{enumerável} (sobre
$d$) de \omterm{def:b2:structures:countable}{conjuntos enumeráveis}, logo é \omterm{def:b2:structures:countable}{enumerável}, e é denso:
$C(\intcc{0}{1})$ é separável.

\textbf{19.} Por linearidade, $\int_0^1 fP = 0$ para todo
polinômio $P$. Escolha polinômios $P_n \to f$ \omterm{def:b2:funcseq:def}{uniformemente}
(Weierstrass):
\[
\Bigl|\int_0^1 f^2\Bigr|
= \Bigl|\int_0^1 f\,(f - P_n)\Bigr|
\leq \norm f_\infty\,\norm{f - P_n}_\infty
\longrightarrow 0 ,
\]
logo $\int_0^1 f^2 = 0$. Se $f(x_0) \neq 0$, a \omterm{def:b2:metric:continuity}{continuidade} dá
$f^2 \geq c > 0$ num subintervalo, contradizendo a integral
nula: $f = 0$. Por consequência, duas funções \omterm{def:b2:metric:continuity}{contínuas} com
os mesmos momentos $\int f t^n$ coincidem.

\textbf{20.} Com $p_{n,k}(x) = \binom nk x^k(1-x)^{n-k}$ e
as convenções $p_{n-1,-1} = p_{n-1,n} = 0$, a regra do produto
e $k\binom nk = n\binom{n-1}{k-1}$, $(n-k)\binom nk =
n\binom{n-1}{k}$ dão
\[
p_{n,k}'(x) = n\bigl(p_{n-1,k-1}(x) - p_{n-1,k}(x)\bigr) .
\]
Somando contra $f(k/n)$ e deslocando o índice na primeira
soma (\omterm{thm:b2:series:abel}{transformação de Abel}):
\[
(B_nf)'(x) = n\sum_{j=0}^{n-1}\Bigl(f\Bigl(\frac{j+1}n\Bigr)
- f\Bigl(\frac jn\Bigr)\Bigr)p_{n-1,j}(x) .
\]

\textbf{21.} Pelo teorema do valor médio, $f(\frac{j+1}n) -
f(\frac jn) = \frac1n f'(\xi_j)$ com $\xi_j \in
\intoo{j/n}{(j+1)/n}$, logo $(B_nf)'(x) = \sum_j
f'(\xi_j)\,p_{n-1,j}(x)$. O nó $\frac{j}{n-1}$ também está em
$\intcc{j/n}{(j+1)/n}$ (as duas desigualdades se reduzem a $j \leq
n-1$), logo $\abs{\xi_j - \frac j{n-1}} \leq \frac1n$ e
\[
\bigl|(B_nf)'(x) - B_{n-1}(f')(x)\bigr|
\leq \sum_j\Bigl|f'(\xi_j) -
f'\Bigl(\frac{j}{n-1}\Bigr)\Bigr| p_{n-1,j}(x)
\leq \omega_{f'}\Bigl(\frac1n\Bigr) \longrightarrow 0
\]
\omterm{def:b2:funcseq:def}{uniformemente}. Como $B_{n-1}(f') \to f'$ \omterm{def:b2:funcseq:def}{uniformemente}
(\cref{thm:b2:funcseq:weierstrass} aplicado à \omterm{def:b2:metric:continuity}{contínua}
$f'$), a desigualdade triangular dá $(B_nf)' \to f'$
\omterm{def:b2:funcseq:def}{uniformemente}. Os polinômios $P_n = B_nf$ convergem então para $f$
no sentido $C^1$.

\textbf{22.} Taylor--Lagrange em $x$: $f(\frac kn) - f(x) =
f'(x)(\frac kn - x) + \frac{f''(\xi_k)}2(\frac kn - x)^2$.
Somando contra $p_k$, o termo linear morre (questão 2):
\[
\abs{B_nf(x) - f(x)}
\leq \frac{\norm{f''}_\infty}{2}\sum_k\Bigl(\frac kn -
x\Bigr)^2p_k
= \frac{\norm{f''}_\infty}{2}\cdot\frac{x(1-x)}{n}
\leq \frac{\norm{f''}_\infty}{8n} .
\]

\textbf{23.} Mais duas derivações de $(x+y)^n$ dão os
momentos fatoriais, com $n_{(j)} = n(n-1)\cdots(n-j+1)$:
\[
\sum_k k_{(j)}\,p_k = n_{(j)}\,x^j \qquad (j = 3, 4),
\]
e $k^3 = k_{(3)} + 3k_{(2)} + k$, $k^4 = k_{(4)} + 6k_{(3)} +
7k_{(2)} + k$ os convertem em momentos de potências:
\[
\sum_k k^3p_k = n_{(3)}x^3 + 3n_{(2)}x^2 + nx,
\qquad
\sum_k k^4p_k = n_{(4)}x^4 + 6n_{(3)}x^3 + 7n_{(2)}x^2 + nx .
\]
Desenvolvendo $(k - nx)^4$ e reunindo (um cálculo paciente mas puramente
mecânico com os quatro momentos de potências):
\[
\sum_k(k-nx)^4p_k = nx(1-x)\bigl(1 + 3(n-2)x(1-x)\bigr) .
\]
Com $x(1-x) \leq \frac14$: o membro da direita é no máximo
$\frac n4\bigl(1 + \frac{3n}4\bigr) = \frac{3n^2}{16} +
\frac n4 \leq n^2$ para $n \geq 1$.

\textbf{24.} A forma de Peano de Taylor em $x$ define $\eta(t) =
\frac{f(t) - f(x) - f'(x)(t-x) - \frac12f''(x)(t-x)^2}
{(t-x)^2}$ para $t \neq x$, $\eta(x) = 0$: por Taylor--Lagrange
$\eta(t) = \frac12\bigl(f''(\xi) - f''(x)\bigr)$ para algum $\xi$
entre $t$ e $x$, logo $\abs\eta \leq \norm{f''}_\infty$ e
$\eta(t) \to 0$ quando $t \to x$ (\omterm{def:b2:metric:continuity}{continuidade} de $f''$). Somando
o desenvolvimento contra $p_k$ e usando as questões 2--3:
\[
n\bigl(B_nf(x) - f(x)\bigr)
= \frac{x(1-x)}{2}f''(x)
+ n\sum_k\eta\Bigl(\frac kn\Bigr)\Bigl(\frac kn -
x\Bigr)^2p_k .
\]
Dado $\varepsilon$, escolha $\delta$ com $\abs\eta \leq
\varepsilon$ em $\abs{t - x}\leq\delta$. Parte próxima: no máximo
$\varepsilon\,n\cdot\frac{x(1-x)}n \leq \varepsilon$. Parte distante:
com $C = \norm{f''}_\infty$ e a questão 23,
\[
n\,C\sum_{\abs{k/n-x}>\delta}\Bigl(\frac kn -
x\Bigr)^2p_k
\leq \frac{nC}{\delta^2}\sum_k\Bigl(\frac kn -
x\Bigr)^4p_k
= \frac{nC}{\delta^2 n^4}\sum_k(k-nx)^4p_k
\leq \frac{C}{\delta^2 n} \longrightarrow 0 .
\]
Logo $n(B_nf(x) - f(x)) \to \frac{x(1-x)}2f''(x)$ ---
o teorema de Voronovskaya. Para $f = e_2$ isso é exato em todo
$n$ (\cref{exo:b2:funcseq:7}): o teto $\frac1n$ é real.

\textbf{25.} (i) A positividade transformou desigualdades \omterm{def:b2:funcseq:def}{pontuais} em
desigualdades de operadores: ela deu a cota da \omterm{def:b2:nvs:norm}{norma}, Jensen,
Chebyshev e todo o Korovkin --- só a linearidade nada demonstra
aqui. (ii) A \omterm{def:b2:metric:compact}{compacidade} entrou por Heine (questões
6 e 13), pela limitação de $f$ e pelo próprio fato de a \omterm{def:b2:nvs:norm}{norma}
$\norm\cdot_\infty$ ser finita. (iii) Três funções de teste
bastam porque a positividade reduz tudo a controlar
$L_n$ na única família $(s-x)^2 = e_2 - 2xe_1 + x^2e_0$,
cujo espaço \omterm{def:b2:structures:generated}{gerado} é o de $e_0, e_1, e_2$. (iv) Bernstein converge
à taxa honesta $\omega(n^{-1/2})$ para toda $f$ \omterm{def:b2:metric:continuity}{contínua}
(ótima, questão 11) mas satura em $\frac1n$ para $f$ suave
(questão 24); o capítulo de Fourier roda o mesmo programa para
funções periódicas com o núcleo de Fejér --- outro operador
positivo com as mesmas virtudes e a mesma modéstia.
\end{solution}
